\section{Conclusion}

This thesis presented PatchOps, a repository-aware, safety-gated, abstention-aware agent for GitHub Actions CI
failures. The project began from a practical goal: diagnose a failed CI run, identify an appropriate repair
action, generate a candidate fix when safe, validate it through CI execution, and avoid harmful automation when
the evidence is weak. The final contribution is more specific and more defensible than simply ``an AI that fixes
CI.'' PatchOps is an AI repair system that treats \emph{when not to fix} as a first-class research problem.

The thesis deliberately avoids claiming to be the first CI-verified repair system or to outperform the strongest
repair baselines on raw Pass@1. Prior CI repair work already establishes repository-level CI-validated repair as
a benchmarked task. PatchOps instead contributes a safety-and-abstention view of the problem, supported by a
working system and a sequence of measured experiments.

\section{Summary of Contributions}

The first contribution is the PatchOps architecture itself: a staged pipeline for CI diagnosis, repository
context retrieval, root-cause reasoning, action planning, patch generation, safety gating, validation, and
decision routing. The system distinguishes between patching, previewing, rerunning, and escalating rather than
treating every failure as a patch-generation task.

The second contribution is an empirical decomposition of CI repair failure modes. The experiments show that
classification becomes reasonably tractable, while fault localization and semantic repair remain difficult.
Across lexical retrieval, LLM hybrid ranking, semantic embedding retrieval, and bounded agentic exploration,
localization plateaus below the target bar. Full-repository retrieval improves patch applicability, but does not
make gated repair achieve CI Pass@1 on the measured repair subset. This isolates the remaining bottleneck as
semantic repository reasoning and repair correctness rather than merely missing files or invalid diffs.

The third contribution is a measured abstention and safety stack. The learned classifier provides a local,
reproducible, low-cost diagnosis stage and a useful confidence signal for selective prediction. Calibration
makes that confidence more usable, and Phase 15 turns the low-risk act-versus-hold decision into a learned,
calibrated gate that replaces the earlier hand-set threshold. The reward-hacking detector asks whether a
CI-green outcome is honest, and the patch-plausibility gate asks whether a patch is likely to be on target
before it is treated as an automatic PR candidate. Together, these components turn abstention from a prompt
instruction into an engineered system behavior.

The fourth contribution is claim discipline. The thesis separates live demonstrations from benchmark evidence,
separates CI validation from semantic correctness, and reports low or negative results rather than hiding them.
That honesty is part of the scientific value of the work: it makes clear where CI repair agents are useful
today and where they remain unreliable.

\section{Answers to the Research Questions}

RQ1 asked whether repository and workflow context improve diagnosis and localization. The results show that
workflow context modestly improves diagnosis, while deeper repository context is more important for file
candidate discovery than for category classification. However, repository retrieval alone does not solve
localization.

RQ2 asked whether PatchOps patches apply and improve CI repair success compared with simpler baselines. The
answer is mixed and intentionally conservative. Full-repository retrieval improves patch applicability among
attempted gated patches, but gated CI Pass@1 remains low on the measured reproducible subset. The raw ungated
baseline can produce some green outcomes, but without the same safety guarantees.

RQ3 asked whether the safety-gated decision policy reduces unsafe or inappropriate repair attempts. The results
support this as the strongest thesis claim. PatchOps explicitly abstains, reruns, previews, or escalates when
automatic patching is not justified, and the later gates add protection against reward hacking and off-target
patches. Phase 15 further strengthens the decision-policy result by showing that the low-risk act-versus-hold
boundary can be learned and calibrated: it matches-or-beats the hand-tuned threshold baseline on the expanded
Safety Set and recovers held-back confident low-risk actions without over-acting on that lane.

RQ4 asked whether PatchOps can distinguish repairable failures from transient or flaky cases. The labeled
safety set shows useful flaky handling and abstention behavior, but the sample remains limited. The conclusion
is that PatchOps has a working decision mechanism for rerun-versus-repair, not that flaky classification is
fully solved.

RQ5 asked about cost and practicality. The measured pipeline remains low-cost in the main offline experiments,
and the learned classifier reduces paid-model dependence for classification. Agentic exploration, however,
adds significant latency without a global localization gain in the measured configuration. This supports the
choice of a conservative staged pipeline with selective use of expensive reasoning.

\section{Future Work}

The most important future direction is semantic repair. PatchOps can now make patches apply more cleanly, but
applicability is not correctness. Future systems should integrate stronger program analysis, differential
testing, invariant checking, and richer execution feedback so that the repair loop can reason about behavior
rather than only about CI outcomes and file overlap.

A second direction is stronger verification. The current reward-hack and plausibility gates are useful but
partial. Future work should add semantic-equivalence checks, mutation-based validation, test-impact analysis,
and stronger patch-correctness oracles. This would make ``CI-green'' one signal among several rather than the
dominant proxy.

A third direction is improved repository navigation. The bounded agentic localizer shows that exploration can
find files missed by retrieval in individual cases, but it is not reliable enough yet. Stronger reasoning
models, graph-guided tools, AST and import-following support, and forced confidence-based exits may make
agentic localization more useful. Future studies should measure these methods against the same localization
and latency baselines rather than assuming agentic navigation is automatically better.

A fourth direction is broader reproducibility. The learned classifier is already local and reproducible, but
the LLM-dependent stages should be evaluated across stronger commercial models and open-weight models with
repeated runs. This would clarify which conclusions are architecture-level and which depend on a particular
model tier.

A fifth direction is larger and more realistic safety evaluation. The current safety set is useful, but real
enterprise CI failures include secrets, protected deployments, permissions, release workflows, flaky
infrastructure, and organization-specific rules. Although the Safety Set has been expanded to support the Phase
15 learned gate, it remains synthetic-heavy. A larger real-derived labeled safety benchmark would strengthen the
abstention claim and make the safety policy more transferable.

A sixth direction is CI-system coverage. PatchOps currently targets GitHub Actions and a restricted local
validation harness. Extending the approach to richer GitHub Actions features, and eventually to GitLab CI,
Jenkins, or other CI systems, would test how much of the architecture transfers beyond the current setting.

\section{Final Remarks}

PatchOps does not solve CI repair, and the thesis does not need it to. The project shows something more useful
for the current state of LLM-based software agents: a CI repair agent must know when not to act. The experiments
demonstrate that localization remains difficult, semantic repair remains unresolved, CI-green is only a partial
oracle, and model self-abstention is not reliable enough by itself. In response, PatchOps builds abstention,
a learned calibrated decision gate, safety gates, plausibility checks, and validation into the repair loop.

The final lesson is therefore practical as well as scientific. Future CI agents should not be judged only by
how often they generate patches. They should be judged by whether they choose the right action under
uncertainty, whether their green results are honest, whether they expose confidence and risk, and whether they
decline to act when the evidence does not justify automation. PatchOps is a step toward that kind of repair
agent: not an overconfident auto-fixer, but a cautious, evidence-aware collaborator inside the CI workflow.
